\scanpage{176}
% archive foliation: 175
\deleted{$\unsure{0}$}
\diagram{\begin{tikzcd}
\mathrm{Prépic}'(X/S)\arrow[r]&\mathrm{Prépic}'(X'/S')\arrow[r,shift left=.5ex]\arrow[r,shift right=.5ex]&\mathrm{Prépic}'(X''/S'')
\end{tikzcd}}
[où $S''=S'\times_S S'$] est exacte \deleted{\unsure{i.e. le} \illegible{2} \unsure{dans}}
\deleted{$\underline{\underline{\mathrm{Prépic}}}_{X/S}$}

\deleted{\unsure{Soit}} En effet, le point (i) est encore une conséquence immédiate du

\uline{Corollaire 1} dans les conditions de Prop. 1)
\margin{($X$ supposé \unsure{q.c. qs} sur $S$)}

[Corollaire 1, soit $S'\to S$ une extension fid. plate \inserted{et q. cpte} de la base. Soit $\underline E$ faisceau \deleted{inversible}\inserted{loc. libre} sur $X$. Pour qu'il soit isomorphe à l'image inverse d'un faisceau loc libre \deleted{[\unsure{représenté}]} sur $S$, il faut que $\underline E'$ soit isomorphe à l'image inverse d'un faisceau loc libre sur $S$.

En effet, il faut exprimer que $f_*(\underline E)$ est loc. libre et \deleted{$\unsure{\xi}$} $f^*f_*(\underline E)-\underline E$ (appl \unsure{naturelle|canonique}) un isom. Or
\[
f_*(\underline E)'=f_*(\underline E')
\]
\margin{extension plate}
on est ramené à des assertions élémentaires, de stabilité de propriétés par extensions fid. plates (chap IV)

